\scanpage{92}\mspage{36}
% L1-3
cofibrés, puisque tout ce qu'on \unsure{pourra|pourrait} en dire revient à ce que je viens d'expliciter, pour $X^{\unsure{\mathrm o|\circ}}$ au lieu de $X$.

% L4-6
Dans la direction de 1\textsuperscript{o}, les arguments précédents à coups de $\Phi_0,p$ nous donnent déjà un renseignement utile.
% L7-11
Mais il sera plus commode à présent de laisser tomber $\underline{\mathrm{Fib}}_0\,S$ ou plutôt, de le remplacer par $\underline{\mathrm{Fib}}\,S$ (ce qui est licite, on l'a dit dans 1\textsuperscript{o}), de sorte qu'on a maintenant un diagramme commutatif (traits pleins)
% L12-14
\diagram{\begin{tikzcd}[ampersand replacement=\&,column sep=small]
 \underline{\mathrm{Fibsc}}\,S\arrow[dr,"j"']\&\&\underline{\mathrm{Fibsc}}_0\,S\arrow[ll,hook,"i"']\arrow[dl,"k"']\\
 \&\underline{\mathrm{Fib}}\,S\arrow[ur,dashed,bend right=38,"{\Phi_0=\unsure{\varphi|\psi}}"']\&
\end{tikzcd}}
,
% L15-18
où $i$ est une inclusion (strictement) pleinement fidèle, $j$ (donc aussi $k$) fidèle. Enfin, on a un foncteur canonique $\varphi(=\Phi_0)$ \deleted{\unsure{tel que}} (en pointillés). Les localis\deleted{\unsure{ations}}eurs sur la ligne du haut sont les localiseurs fins, (pour $\underline{\mathrm{Fibsc}}_0\,S$
% L19-23
c'est pareil que la $W$-équivalence grossière), sur $\underline{\mathrm{Fib}}\,S$ par contre on prend la $W$-équivalence grossière. (On sait que le localisé est le $=$, à équivalence près, que si on se restreint à la sous-catégorie pleine $\underline{\mathrm{Fib}}_0\,S$, ce qui est assez remarquable.)
% archive sheet number: 91
